O avanço tecnológico dos Veículos Aéreos Não Tripulados (VANTs) ampliou suas possibilidades de aplicação, o que permitiu superar as limitações de mobilidade dos robôs terrestres em tarefas como inspeção, agricultura e transporte de objetos para locais inacessíveis \cite{carvajal2024}. Com essa evolução, os VANTs deixaram de funcionar só como plataformas de observação e passaram a ser capazes de executar tarefas autônomas que envolvem interação física com o ambiente, como missões de \textit{pick-and-place}, focadas na captura e posicionamento de objetos \cite{abuzayed2020, cocuzza2021}. A execução autônoma dessas tarefas exige o desenvolvimento de algoritmos de controle que permitam que o robô opere com segurança \cite{carvajal2024}. Além disso, em ambientes nos quais os sistemas de navegação global (GPS) não oferecem a precisão necessária para o alinhamento de uma garra robótica, o servocontrole visual, também chamado de servovisão, surge como a solução mais viável, utilizando o retorno visual de câmeras embarcadas para orientar o VANT \cite{shang2016}.

A captura autônoma de um alvo não cooperativo, como, por exemplo, armadilhas terrestres de monitoramento entomológico sem marcadores visuais fiduciais, impõe desafios significativos aos métodos clássicos de percepção e à estabilidade da aeronave \cite{zhao2019}. A motivação para investigar esse problema está relacionada à vulnerabilidade dos sistemas tradicionais de visão quando operam em ambientes externos e não estruturados \cite{zhao2019}. No que diz respeito à percepção, algoritmos baseados exclusivamente em modelos de espaço de cores, como HSV, são altamente sensíveis às variações de iluminação. Da mesma forma, técnicas clássicas de extração de características geométricas, como SIFT, SURF e HOG, combinadas com estratégias de varredura por janelas deslizantes, exigem esforço computacional redundante e apresentam dificuldades para modelar a diversidade dos planos de fundo das imagens \cite{zhao2019}. Como consequência, essas abordagens clássicas frequentemente sofrem com a chamada lacuna semântica (\textit{semantic gap}), demonstrando incapacidade de abstrair o contexto da cena, o que leva o sistema a confundir a ausência de luz em um objeto escuro com sombras ou ruídos do próprio ambiente \cite{zhao2019}.

Em relação ao controle, a servovisão sofre com atrasos, uma vez que o processamento de imagens é computacionalmente intensivo e opera em frequências significativamente inferiores às dos sensores inerciais de voo. Isso resulta em um período de amostragem elevado e, consequentemente, gera intervalos sem retorno visual \cite{shang2016}. Além disso, em missões de captura aérea baseadas em mecanismos acoplados sob o \textit{frame} da aeronave, o sistema visual precisa fornecer não apenas o posicionamento espacial do alvo, mas também referências para o alinhamento preciso da orientação de guinada (\textit{yaw}) e para a compensação do deslocamento físico entre a câmera e o atuador durante a descida \cite{cocuzza2021, ladig2021}.

Diante da necessidade de superar essas limitações tecnológicas, formula-se a seguinte questão de pesquisa: como integrar algoritmos de detecção baseados em aprendizado profundo a estratégias de controle dinâmico para viabilizar a captura autônoma de um alvo não cooperativo por um VANT, mesmo diante de variações nas condições de iluminação e de atrasos no processamento de imagens, garantindo o alinhamento preciso das garras robóticas durante a operação?

Para contextualizar o problema investigado, este trabalho integra o projeto VTS (\textit{Vector Tracking System}), financiado pela Fundação de Apoio à Pesquisa do Distrito Federal (FAPDF) e desenvolvido pelo IDP em parceria com a Universidade de Brasília (UnB) e a Fiocruz de Recife. O projeto tem como objetivo principal desenvolver soluções tecnológicas para a vigilância e o mapeamento automatizado do mosquito \textit{Aedes aegypti}. Para isso, contempla a integração de diferentes tecnologias, incluindo armadilhas inteligentes, sistemas de telemetria baseados em LoRaWAN, modelos de classificação biológica e aeronaves não tripuladas autônomas.

Nesse contexto multidisciplinar, esta pesquisa concentra-se no desenvolvimento de um sistema de transporte e captura aérea de armadilhas. Enquanto outras frentes do projeto se dedicam ao desenvolvimento da eletrônica embarcada e dos sensores que compõem esses dispositivos, este trabalho aborda os mecanismos de percepção visual e controle de voo necessários para que um VANT seja capaz de identificar a estrutura da armadilha, realizar os ajustes de posicionamento e efetuar o acoplamento físico de maneira autônoma. Dessa forma, busca-se contribuir para a automatização dos processos de captura e recolhimento das armadilhas, especialmente em regiões remotas ou de difícil acesso.

Para enfrentar os desafios relacionados à detecção, a literatura tem se concentrado cada vez mais em modelos de aprendizado profundo (\textit{Deep Learning}) \cite{amjoud2023}. Nesse contexto, a família YOLO (\textit{You Only Look Once}) de redes neurais convolucionais destaca-se por realizar a detecção de objetos em uma única passagem pela rede, proporcionando um equilíbrio entre precisão e velocidade de processamento \cite{afifah2025, amjoud2023, adegun2024}. Considerando a evolução dessas arquiteturas e a necessidade de atender às restrições computacionais de sistemas embarcados em VANTs, torna-se relevante investigar o desempenho de diferentes gerações desses modelos. Assim, este trabalho propõe uma avaliação comparativa das arquiteturas YOLOv8n, YOLOv10n, YOLO11n e YOLO26n, buscando identificar aquela que oferece o melhor equilíbrio entre precisão e custo computacional para a detecção embarcada de armadilhas.

Embora as redes neurais ofereçam alta precisão na extração de características semânticas, a percepção visual em VANTs é frequentemente afetada por ruídos dinâmicos. A detecção de objetos com VANTs está sujeita a falhas momentâneas e mudanças abruptas nas caixas delimitadoras devido às vibrações que ocorrem durante o voo \cite{barisic2019}. Para atenuar essas falhas e as oscilações ruidosas das redes neurais aplicadas a drones, o filtro de Kalman tem sido utilizado como uma solução eficaz \cite{barisic2019}. Essa técnica de estimação permite obter uma aproximação contínua e confiável não apenas da posição do alvo, mas também de sua velocidade cinemática, compensando as deficiências geradas durante períodos de perda do rastreamento visual \cite{barisic2019}.

Além da filtragem das estimativas visuais obtidas pelo detector, o referencial espacial adotado para a navegação representa outro desafio. Isso decorre das limitações de implementar o laço de controle do VANT exclusivamente no plano bidimensional da imagem (2D). Quando a navegação se baseia apenas em coordenadas de pixels, o sistema torna-se suscetível à ambiguidade de profundidade e sensível a incertezas geométricas, uma vez que variações de altitude alteram a escala do objeto na imagem capturada pela câmera. Nesse contexto, existem vantagens significativas na aplicação do modelo analítico de câmera \textit{Pinhole} \cite{thuilot2002, fagundes2025}. O uso de matrizes de transformação geométrica permite converter diretamente as coordenadas dos pixels em um espaço cartesiano métrico tridimensional (3D), reduzindo a sensibilidade a incertezas durante a fase de aproximação e proporcionando uma dinâmica de voo substancialmente mais estável \cite{fagundes2025}.

Com base nesse cenário, a hipótese central deste trabalho é que a integração de um algoritmo de detecção de objetos, responsável por identificar e delimitar semanticamente um alvo não cooperativo, com um filtro de Kalman e uma transformação espacial de duas para três dimensões (2D para 3D), fundamentada no modelo de câmera \textit{Pinhole}, seja capaz de reduzir os efeitos de ruídos nas detecções visuais e fornecer estimativas de posição mais estáveis e precisas. Essas informações, por sua vez, devem viabilizar o controle de voo necessário à execução autônoma da captura de objetos por VANTs utilizando técnicas de servovisão.

Adicionalmente, formula-se a hipótese de que o emprego da velocidade estimada e filtrada pelo filtro de Kalman como variável de entrada de um controlador proporcional-derivativo (PD) em cascata permita compensar os efeitos dos atrasos inerentes à amostragem e ao processamento das imagens, além de reduzir a ultrapassagem do ponto de referência (\textit{overshoot}). Essa estratégia busca evitar a dependência do termo integral, presente em controladores PID tradicionais, cuja contribuição pode agravar as respostas transitórias diante de atrasos e variações na dinâmica do sistema \cite{shang2016, jiang2021, fagundes2025}.

Diante do exposto, o objetivo geral desta monografia é desenvolver e validar, em ambiente simulado, um sistema autônomo de captura de objetos por VANTs, utilizando técnicas de servovisão (\textit{Visual Servoing}) para detectar e localizar um alvo não cooperativo, realizar a aproximação e controlar o posicionamento da aeronave de modo a garantir o alinhamento necessário à execução da captura. Para alcançar o objetivo geral, foram definidos os seguintes objetivos específicos:

\begin{enumerate}

	\item Treinar e avaliar modelos de diferentes gerações da família YOLO nas três fases operacionais do voo (aproximação macroscópica em trajetória de funil, alinhamento angular de guinada e alinhamento fino de posição com compensação de \textit{offset} em baixa altitude), buscando atingir precisão média superior a 90\% na detecção do alvo;

	\item Desenvolver e avaliar um sistema de processamento de imagens capaz de operar a uma taxa mínima de 20 quadros por segundo (fps), fornecendo informações visuais em tempo hábil para o controle autônomo do VANT;

	\item Aplicar e avaliar um filtro de Kalman para atenuar os ruídos nas estimativas de posição e orientação angular da armadilha, mantendo a predição de sua trajetória por até 4,0 segundos durante interrupções temporárias da detecção visual;

	\item Estabelecer uma estratégia de controle visual com mecanismo de histerese para manter a armadilha no campo de visão da câmera durante a descida do VANT, utilizando a transformação de coordenadas da imagem em coordenadas espaciais para orientar os ajustes de posicionamento;

	\item Desenvolver e avaliar uma estratégia de alinhamento angular autônomo do VANT em relação à armadilha, buscando limitar o erro de guinada (\textit{yaw}) a menos de $1{,}0^\circ$ em relação à orientação de referência;

	\item Determinar e avaliar a precisão do alinhamento posicional do VANT em relação ao ponto de captura, buscando obter raiz do erro quadrático médio (RMSE) inferior a 8 mm;

	\item Validar o desempenho do sistema por meio de simulações de missões autônomas completas, abrangendo a aproximação, os alinhamentos angular e posicional, o acoplamento e a descida vertical cega, buscando alcançar uma taxa de sucesso superior a 60\% nas tentativas de captura.

\end{enumerate}

A metodologia desta pesquisa caracteriza-se como um estudo de modelagem e validação em ambiente integralmente simulado. Inicialmente, são coletadas imagens para a construção de conjuntos de dados (\textit{datasets}) específicos, utilizados no treinamento supervisionado e na avaliação de diferentes arquiteturas da família YOLO nas três fases operacionais do voo: aproximação macroscópica, identificação das alças para alinhamento angular de guinada e alinhamento posicional fino em altitude muito baixa. Os modelos são comparados quanto ao desempenho na detecção do alvo, à precisão geométrica e à latência computacional. Em seguida, aplica-se um filtro de Kalman para estimar os estados do sistema e atenuar os ruídos das detecções visuais. As coordenadas bidimensionais da imagem são então convertidas em estimativas métricas de posição no espaço tridimensional, por meio do modelo de câmera \textit{Pinhole}, considerando a compensação do deslocamento físico entre a lente e o ponto de referência da captura.

Na etapa seguinte, desenvolve-se uma estratégia de controle de voo baseada em controladores proporcionais-derivativos (PD) em cascata, que utiliza as estimativas filtradas para orientar a aproximação translacional, o alinhamento angular de guinada e o posicionamento fino da interface de captura das garras em relação ao alvo. Por fim, os módulos são integrados por meio do ROS 2, responsável pela comunicação dos dados provenientes da câmera simulada no Gazebo Sim. A dinâmica de voo é gerenciada pelo piloto automático ArduPilot em arquitetura \textit{Software-in-the-Loop} (SITL), com transmissão dos comandos de controle e dos dados de telemetria pelo protocolo MAVLink. Essa integração permite avaliar o desempenho do sistema em simulações de missões autônomas, abrangendo as etapas de detecção, aproximação, alinhamento e captura do alvo.

Por fim, o restante deste trabalho está organizado em quatro capítulos. O Capítulo~2 apresenta a revisão de literatura, contemplando a fundamentação teórica sobre VANTs, visão computacional, servovisão e manipulação aérea, além da análise comparativa dos trabalhos correlatos. O Capítulo~3 detalha a metodologia proposta, incluindo o ambiente de simulação, a geração dos conjuntos de dados, a filtragem estocástica, o controle em cascata e a máquina de estados da missão. O Capítulo~4 expõe e discute os resultados experimentais obtidos nas campanhas de validação e nos ensaios de Monte Carlo. Por fim, o Capítulo~5 sintetiza as conclusões da pesquisa, suas limitações e as perspectivas para trabalhos futuros.
